\textbf{Inherited tolerance, prior injury and local repair answer different questions.}
The density fields and curves in panels A and B use the completed native outer-search package (2432 unique updates, four candidates,40 stages and 24 complete paths). The compact parent and offspring are 93e6148d32bbb0862163 and 8d60f7d250e8e9f2a25b; the multiring pair is d02e7f60565f83066e5d and a6a821dffd4f4046c97a. Each offspring is one compound kernel/initial-body mutation of its parent. At life 0, all four have the same total mass 71.248610437 and energy 56.998888349, with E=0.8 rho and common resource fields; their later bodies and inventories differ. All domains are 96 by 96 and periodic. There is one common forcing/random tape, one child per parent and one fixed+x injury scenario, not independent replication or a pure shape intervention. Physical rules are fixed within each life; learn=false, open uptake and fresh-brake at natural expiry are retained.

\textbf{Design and clocks.} Each candidate forms over updates 1--128 with T00, which disables both transport capability and its fee. Its full state 128 branches into three histories over 129--160. N (naive) uses T11 paid P/F transport without a prior cut. C (experienced) also uses T11, with 10\%+x cuts immediately before updates 129 and 145. The T00 history has no prior cut and no transport; it is a policy baseline, not a pure cost control. Every actual 160 receiver branches to no additional cut (none) or a 20\%+x cut immediately before 161 (x), then continues to 224 with its history's transport setting. Each candidate therefore contributes 128+3 x 32+6 x 64=608 unique updates. Displaying the shared formation/acquisition prefixes on complete paths does not count or execute them again. Body-time is 0.05 times the sum of old body mass after any declared pulse and before each update. The figure's full-cycle differences include the whole 0--224 history and do not substitute future-only 64-step totals.

\textbf{A: inherited survival versus the value of experience.} Original-coordinate density fields at 160, immediately after the test cut before 161,184 and 224 use the same inclusive crop x 28--67,y 28--63 for all compact rows. N/C are both represented for parent and offspring. Coordinates are not translated or rotated; y increases upward. White denotes zero density and teal denotes 1 M/cell on a shared linear scale. The native zero annotations mark actual all-zero fields, not unobserved states. The compact parent's challenged N/C/T00 paths reach exact zero at 178/183/177 and stay zero through 224. Its offspring survives all six futures through 224, including all three challenged histories. The six native mass polylines show all compact challenged histories over 0--224: lighter colours denote the parent, darker colours the offspring; blue=N, orange=C and purple=T00. All 225 states including life 0 are supplied; the display geometry alone is simplified at 0.22 authoring-screen pixels before integer-pointer placement. Axes start at zero and the vertical guides mark 128 and 161. The image immediately after the pulse is a separate saved receipt; ordinary trajectory points are post-update states.

For naive/x, offspring minus parent full-cycle body-time is+213.152894467, and the offspring ends at M 48.232800969/E 33.138175905 versus exact-zero parent. The actual test-cut mass is 12.183807878 versus 8.766334622: survival is not explained by a smaller absolute cut. However, within the offspring, C-minus-N full-cycle body-time is-24.084535875 without a test cut and-46.595738718 with a cut; corresponding endpoint M differences are-13.520030444/-5.536194159 and E differences-8.269522115/-11.898291854. The naive challenged body has two peaks while the experienced challenged body develops open/ring-like contours; these shape differences do not establish a higher return. C/x is also below the T00/x baseline in full-cycle body-time(-14.064205300), M(-0.086397780) and E(-1.054764198). All three challenged offspring lose net M in the last 32 updates: N-10.056520173,C-1.226291635,T00-0.650326685, with 31/24/22 negative steps. Survival to 224 is not stable restoration or long-term safety.

\textbf{B: a failed multiring mutation and a positive interaction without net benefit.} The left images show the multiring N/C challenged parent and offspring at 160 and 224 in the same inclusive x 14--81,y 14--77 crop, with y increasing upward. The colour scale is the same as A, but spatial magnification differs; displayed areas must not be compared as masses. The offspring grows more slowly during formation and remains inferior in all six history-by-challenge comparisons: endpoint M, endpoint E and full-cycle body-time are all lower than the matched parent. These six comparisons are N/none,N/x,C/none,C/x,T00/none,T00/x; none is omitted as a successful exception. The naive/x offspring ends at M 348.399796679/E 203.765446202 versus parent 553.768566916/289.653255642, with full-cycle body-time lower by 927.498240914. Its actual test-cut mass is 54.165575426 versus parent 79.605977139. Both multiring parent and offspring gain mass at every update of the final 32 in all six futures, unlike the compact challenged offspring.

The right native curves plot cumulative full-cycle C-minus-N body-time for each offspring and both futures, from state 0 through 224. The identical formation prefixes give zero differences; differences after the experienced cuts are negative. Blue/orange identify compact none/x and purple/teal multiring none/x. All 225 archived states are supplied to the same time and body-time scales. Multiring endpoint differences are-56.515604586 (none) and-50.218078800 (x). Their subtraction gives the positive interaction  $I_T$ = +6.297525786; it does not make either complete experienced path better than naive. Endpoint M/E are also lower in both cases. Multiring C/x exceeds T00/x in terminal M by 23.703860033 but has lower E(-2.058940748) and complete body-time(-15.880144833), so terminal M alone hides a tradeoff. All original fees, pulse removals and negative intervals remain in the source tables. Zero starvation and no substantial unpaid charges in this package do not identify a sufficient explanation for the mortality or survival differences.

\textbf{Separate rule--initial-state follow-up (A matrix and B annotation).} The added matrix and annotation use the completed 1792-update package, job d7dcd928-7204-4e24-ab92-1bd791394e7e. Four new off-diagonal R/I combinations use 28 new stages; corresponding parental and offspring controls reuse 28 old stages (1792 unique old updates) from the original 2432 package. Those reused updates are not new executions. R0/I0 denotes each parental rule/initial-state package; R1/I1 denotes its offspring. R consists of the two actual kernel arrays and the complete response rule. I consists of density $\rho$, its corresponding energy distribution, founder field and neutral polarity distribution. The intervention is at package level, not a single-gene or pure-shape manipulation. Each new combination uses the same formation 0--128, N/C histories 129--160 and none/x tests 161--224. There are four matched history--challenge paths at each I: N/none, N/x, C/none and C/x. T00 histories in the original six-path comparisons are not added to this follow-up.

The compact matrix lists N/x and C/x exact-zero updates in that order: R0I0 178/183, R0I1 178/184. Both challenged paths of R1I0 and R1I1 retain nonzero mass through 224. These are actual extinction updates, not first saved zero-frame times. This finite-window survival grouping does not remove initial-state effects or rule--initial-state interaction. Under compact R1I0, C-minus-N at x has endpoint mass $+0.686585896$, endpoint energy $-6.279353502$ and complete 0--224 body-time $-40.632697448$; a positive endpoint mass alone does not demonstrate a benefit of prior injury. The original plotted trajectories, images and zero tails remain unchanged.

For multiring bodies, R1-minus-R0 endpoint M, endpoint E and full-cycle T are negative for all four matched history--challenge paths at each of I0 and I1. Rule--initial-state interactions are nonzero: the four body-time interactions are $+23.247941681$, $+15.550813277$, $+22.243867153$ and $+25.883975770$ (N/none, N/x, C/none, C/x). These factorial interactions differ from the original C-by-injury interaction plotted in B. They do not reverse the negative rule effects. All 16 multiring futures in the factorial comparison remain nonzero and gain mass at every final-32 update; the compact late decline does not generalize to them. This follow-up uses the same tape after selection of the original parents and offspring and is not an independent replication or a claim of learned recovery.

\textbf{C: separate evidence for partial local regrowth without shape restoration.} This inset is from the earlier independent excision-transfer package, not a continuation of A/B. Two existing open-T11 state 256 receivers, with their original fields/rules/clock, were cut once by 20\% along+x and followed over 257--384. Two injured continuations contribute 256 original updates; their two uninjured same-clock references reuse 256 unique archived packets without replay. The compact receiver loses all positive generation requests/births by 261 and reaches exact zero at 273, retaining its zero tail. The multiring injured body regains its pre-cut total M at 312 but still ends below its own uninjured reference by 104.616522 M and 53.505709 EU.

The two density fields show the multiring injured and uninjured bodies at 384 in the same inclusive x 10--85,y 10--85 crop with the same density scale and orientation as above. The native curve uses all 129 states 256--384 for mass in the fixed original cut support, defined by removed rho>0, including partly cut cells. It is neither a moving region nor a complete rectangular column mask. The injured cut region starts at 2.549621981, peaks at 36.703056814 at 325, then loses mass continuously over 326--384 to 33.215383281. Its uninjured reference ends at 145.837073833; the injured region is 22.78\% of that value. Cumulative local births 51.423284391, local rule deaths 15.968602588 and net tissue input-4.788920503 close the region's mass increase 30.665761300. This is genuine local generation with net export, not restoration by influx alone, and is a path account rather than a causal mediation decomposition. Current regional tissue contains 30.799173 mass carrying the passive local-birth marker (92.73\%); this is not an independent lineage. Whole-body mass still increases 172.973381 over the post-cut window while this region declines late. Raw and integer-registered density L1 are both 373.060436870 with displacement(0,0): total M recovery does not imply restoration of the missing+x structure. The two experiments have different receivers and clocks and are not combined into a sequential mechanism or a timing effect.

All panels were assembled with visible Impress pointer/keyboard actions. Curves, labels and axes remain native editable objects; only individual unlabelled scalar fields are raster assets extracted from archived arrays. No diagnostic whole figure is pasted, and no model or energy forward step is executed for authoring. Selection into an evolutionary archive is not asserted. The candidate preserves failure, finite-window decline, the strong T00 baseline and the distinction between inherited rules and changing individual state.
